\section{Setup and Lawvere's Diagonal}
\label{sec:setup}

We fix notation, name three predicates on systems, and state Lawvere's theorem in its simplest form. Readers who already know \citet{lawvere1969} or \citet{yanofsky2003} can skip ahead; the predicate definitions in \S\ref{subsec:predicates} are the load-bearing pieces.

\subsection{Systems as functions of two arguments}
\label{subsec:system-function}

A system is a function. We write
\begin{equation}
\System : A \to (A \to Y),
\qquad \text{or curried} \qquad
\System : A \to A \to Y,
\end{equation}
and read the first argument as a self-description (a system prompt, a context, anything that configures behavior) and the second as the input the configured system actually receives. The two-slot shape is what makes self-application well-typed: $\System(a)(a)$ is the value of the system when its self-description and its input coincide. That diagonal is where everything in this paper happens.

For language models the natural choice is $A = \Prompt$ and $Y = \Token$. Nothing in the framework cares about that choice, though; any deterministic system with a description--input split fits.

\begin{definition}[Diagonal]
\label{def:diag}
The \emph{diagonal} of $\System$ is $\diag(\System) : A \to Y$, $\diag(\System)(a) := \System(a)(a)$.
\end{definition}

\subsection{Three predicates}
\label{subsec:predicates}

\begin{definition}[Universal]
\label{def:universal}
$\System$ is \emph{universal}, written $\Universal(\System)$, if it is surjective in the curried sense: every $g : A \to Y$ equals $\System(a)$ for some $a \in A$.
\end{definition}

\begin{definition}[Controlled]
\label{def:controlled}
Fix an endomap $t : Y \to Y$. Read $t$ as the transformation a controller wants to forbid: ``flip safe to unsafe,'' ``negate the claim,'' ``replace \texttt{refuse} with \texttt{comply}.'' The system is \emph{$t$-controlled}, written $\Controlled_t(\System, p)$, when the self-prediction $p : A \to Y$ never lands on a $t$-fixed-point:
\begin{equation}
\forall a \in A, \quad t\bigl(p(a)\bigr) \neq p(a).
\end{equation}
\end{definition}

\begin{definition}[Transparent]
\label{def:transparent}
A self-prediction $p : A \to Y$ is supposed to anticipate, from outside, what $\System$ does on its own description. The system is \emph{$p$-transparent}, written $\Transparent_p(\System)$, if
\begin{equation}
\forall a \in A, \quad \System(a)(a) = p(a).
\end{equation}
\end{definition}

The three predicates pick out three things the AI safety literature wants at once: capability (Universal), constraint (Controlled), and honesty about oneself (Transparent). \Cref{sec:mirror-trilemma} shows that any two are jointly satisfiable but the conjunction of all three is contradictory.

\subsection{Lawvere's theorem}
\label{subsec:lawvere}

\begin{theorem}[\citealp{lawvere1969}]
\label{thm:lawvere}
Let $A$, $Y$ be sets, $f : A \to A \to Y$ surjective. Then every endomap $t : Y \to Y$ has a fixed point.
\end{theorem}

\begin{proof}
Define $g(a) := t(f(a)(a))$. By surjectivity, $f(a_0) = g$ for some $a_0$. Evaluate at $a_0$:
\[
f(a_0)(a_0) = g(a_0) = t\bigl(f(a_0)(a_0)\bigr).\qedhere
\]
\end{proof}

\begin{corollary}[Contrapositive]
\label{cor:contrapositive}
If $t$ has no fixed point, no surjective $f : A \to A \to Y$ exists.
\end{corollary}

\begin{remark}
\label{rem:scope}
The same diagonal subsumes Cantor's theorem ($Y = \Bool$, $t = \neg$), Russell's paradox, G\"odel's first incompleteness, Tarski's undefinability, the halting problem, and Rice's theorem under different choices of $t$ and $Y$; see \citet{yanofsky2003} for the unified categorical treatment. \Cref{sec:landscape} adds eight AI-safety impossibilities to the list.
\end{remark}

\subsection{The constructive witness}
\label{subsec:diagonal-point}

The proof above is constructive. Reading off the data:

\begin{definition}[Lawvere diagonal point]
\label{def:diagonal-point}
Given surjective $f : A \to A \to Y$ and an endomap $t : Y \to Y$, set $g(a) := t(f(a)(a))$, pick any $a_0$ with $f(a_0) = g$, and define
\begin{equation}
y_\star := f(a_0)(a_0).
\end{equation}
$y_\star$ is the \emph{Lawvere diagonal point} of $(f, t)$, and $a_0$ is the \emph{diagonal prompt}.
\end{definition}

For a language model with $\System : \Prompt \to \Prompt \to \Token$, $a_0$ is an actual prompt --- a token sequence --- where the system's self-application equals what the controller forbids. The existence of $a_0$ is what drives every impossibility in this paper. Whether $a_0$ is reachable in practice, and what its cost looks like, is empirical and we say what we can in \Cref{sec:discussion}.

\medskip
With the framework in place, \Cref{sec:mirror-trilemma} states and proves the Mirror Trilemma.
